\documentclass[10pt,letterpaper]{article}

\usepackage[T1]{fontenc}
\usepackage{lmodern}
\usepackage[margin=0.60in]{geometry}
\usepackage{array}
\usepackage{booktabs}
\usepackage{tabularx}
\usepackage{enumitem}
\usepackage[hidelinks]{hyperref}

\setlength{\parindent}{0pt}
\setlength{\parskip}{3pt}
\setlist[itemize]{leftmargin=1.25em,itemsep=1pt,topsep=2pt}
\renewcommand{\arraystretch}{1.12}
\newcommand{\openstatus}{\textbf{Open}}
\newcommand{\partialstatus}{\textbf{Partly supported}}

\begin{document}
\small

\begin{center}
{\Large\bfseries Remote Environmental Survey and Mapping of an Unseen Room}\\[4pt]
Julian Joseph Amaro \quad Project I (247-519-VA) \quad Vanier College\\
November 20, 2026
\end{center}

\section*{1. Change Log}

The submitted proposal remains the frozen baseline for needs N1--N8, watch constraints C1--C5, requirements REQ-RS-01--08, and risks R1--R3. The watch and custom-PCB commitments already appear there; the new identifiers formalize them without changing the intended scope.

\begin{tabularx}{\linewidth}{>{\raggedright\arraybackslash}p{0.16\linewidth}X>{\raggedright\arraybackslash}p{0.16\linewidth}}
\toprule
\textbf{Item} & \textbf{Change and reason} & \textbf{Record} \\
\midrule
REQ-RS-09 & Formalized the proposal's watch-output commitment as a numbered requirement: the exact Waveshare ESP32-S3-Touch-AMOLED-2.06 displays the map, summary, and anomaly indication. & L-02, L-03.1 \\
REQ-RS-10 & Formalized the proposal's custom-hardware/Project II PCB commitment as a numbered final-system requirement. Stage 1 may use development modules, breadboard, perfboard, or protoboard. & L-02, L-03.1 \\
Gas interpretation & Clarified that the BME688 is a relative environmental/anomaly sensor, not selective toxic-gas detection or proof that a room is safe. & L-02--L-05, L-03.1 \\
\bottomrule
\end{tabularx}

A Raspberry Pi Zero 2 W and a 3D-printed tracked tank are candidate components only. Neither is a frozen architecture decision or a verified compatible choice.

\section*{2. Concept Alternatives}

Two approaches differ in where time-sensitive control runs. In A, one computer handles video, mapping, sensing, and motor timing. In B, a separate controller handles motors and encoders while the main computer handles video, mapping, storage, and survey transfer.

\begin{tabularx}{\linewidth}{X>{\centering\arraybackslash}p{0.10\linewidth}>{\centering\arraybackslash}p{0.12\linewidth}>{\centering\arraybackslash}p{0.12\linewidth}}
\toprule
\textbf{Criterion (1--5; higher is better)} & \textbf{Weight} & \textbf{A} & \textbf{B} \\
\midrule
Requirement coverage & 0.35 & 3 & 4 \\
Motor/encoder timing and mapping risk & 0.25 & 2 & 4 \\
Project I buildability and obtainable evidence & 0.20 & 4 & 3 \\
Power and integration simplicity & 0.10 & 4 & 2 \\
Cost headroom & 0.10 & 4 & 3 \\
\midrule
\textbf{Weighted total} & \textbf{1.00} & \textbf{3.15} & \textbf{3.50} \\
\bottomrule
\end{tabularx}

\textbf{Provisional choice: B.} Its separation of motor timing from video/map processing reduces a key risk; the phone-controlled ESP32 rover and separate video PoCs support exploring it. A loses on timing; B loses on simplicity and cost and remains unverified as an integrated system. The five weights and 1--5 rubric were fixed in the dated L-09 record before this finalized comparison. The scores are engineering judgments, not measured performance, and need review after integration tests.

\textbf{Functional blocks for provisional B:}
\begin{center}
\fbox{Phone command} $\longrightarrow$ \fbox{Motion/encoder control} $\longrightarrow$ \fbox{Rover motion}\\[3pt]
\fbox{Camera, range, environment} $\longrightarrow$ \fbox{Video/map processing} $\longrightarrow$ \fbox{Phone view}\\[3pt]
\fbox{Map and tagged readings} $\longrightarrow$ \fbox{Summary transfer} $\longrightarrow$ \fbox{Watch display}
\end{center}

\newpage
\section*{3. Feasibility Case}

Claims begin at C6 because C1--C5 already identify watch constraints in the frozen proposal.

\begin{tabularx}{\linewidth}{>{\bfseries}p{0.055\linewidth}>{\raggedright\arraybackslash}p{0.12\linewidth}X>{\raggedright\arraybackslash}p{0.18\linewidth}>{\raggedright\arraybackslash}p{0.12\linewidth}}
\toprule
ID & Tests & Claim and present evidence & Class / source & Status \\
\midrule
C6 & RS-06, RS-09 & The exact watch ran V1/V2 Wi-Fi scans and a touch list (L-06; photos). V2 built in isolation. A standalone phone-photo upload/display worked per user report (L-08). The local SquareLine3 combination compiled November 17; the user now reports flashing it and displaying a phone-uploaded image within that UI. A rover-generated map, survey records, anomaly, and 5 s delivery remain open. & Measured V1/V2; user-reported SquareLine3 physical test; builds; L-06, L-08, L-10 \cite{waveshare} & \partialstatus \\
C7 & RS-03, RS-04 & The BME688 documents temperature, humidity, pressure, and broad gas response. It can support a relative anomaly channel, but it is not selective or certified toxic-gas detection. & Documented research; L-02, L-03 \cite{bosch} & \partialstatus \\
C8 & RS-04 & Seven BME688 project pages were reviewed for possible data-collection approaches. Their code and data sets were not downloaded or run; no result was reproduced. This review does not establish gas-species identification, which is outside the current relative-anomaly scope. & Page review only; L-05, L-10 & \openstatus \\
C9 & RS-05 & A wheeled project-box rover with three front ToF sensors and encoder-equipped N20 motors was driven by phone Wi-Fi while a live phone map used sensor readings, per user report. Encoder contribution was not demonstrated; no map points were tape-measured. The proposed tracked tank remains untested. & User-reported physical PoC; L-10; tank assumption L-02 & \openstatus \\
C10 & RS-01, RS-02, RS-05 & The Pi Zero 2 W has not been tested. The separate laptop/desktop-to-phone video trial does not establish its video/mapping load, thermal behaviour, motor timing, or power demand. & Assumed Pi capability; L-09, L-10; camera record & \openstatus \\
C11 & RS-01--RS-10 & Phone drive and a ToF-based live map worked together on the wheeled rover, while laptop video and SquareLine3 watch photo display were separate tests, per user report. No complete rover-to-watch survey or integrated acceptance run exists. & User-reported subsystem PoCs; L-10; integrated tests pending & \openstatus \\
C12 & RS-01 & User reports an ESP32/N20 rover drove farther than 2 m by phone Wi-Fi using \texttt{sketch\_oct31b}. The redacted repository copy compiled; calibrated distance record and outside-room operator position are absent. & Reported physical PoC; L-09; dated rover V2 record & \partialstatus \\
C13 & RS-02 & User reports a live laptop/desktop camera feed on a phone; a phone-browser screenshot and Flask/OpenCV source are preserved. A still frame cannot establish motion or delay. Raspberry Pi operation and $\leq$1 s latency are untested. & Reported physical PoC; L-09; Nov. 19 camera record & \partialstatus \\
C14 & RS-03--RS-05 & Python simulations generate three synthetic tagged readings, distinguish 124 from 125 at baseline 100, and score eight wall points across five synthetic runs. All 10 software tests passed. No physical sensor or mapping result follows. & Calculated/simulated; L-09; simulation outputs & \partialstatus \\
C15 & RS-07 & A provisional CAD \$675 parts estimate leaves no verified margin after tax, shipping, and changes. Itemized final receipts are absent. & Calculated estimate; L-03.1 & \openstatus \\
C16 & RS-08 & Separate watch and rover PoC builds passed, but no second computer built the integrated firmware from a clean copy with one command. & Measured builds; L-06, repository build records & \partialstatus \\
C17 & RS-10 & The final custom PCB is a binding plan, not a fabricated or tested board. & Assumed; proposal and L-02 & \openstatus \\
\bottomrule
\end{tabularx}

These PoC results support subsystem feasibility, not complete requirement compliance. Rover distance, its ToF-based live map, and the live feed are user-reported; the screenshot and code cannot establish missing acceptance conditions. Simulated 03--05 outputs are synthetic. No BME680/BME688 bench trial, tape-measured mapping-accuracy trial, Pi video-latency trial, rover-to-watch survey transfer, or full clean-copy integrated build has been documented.

\section*{4. Baselined Requirements}

\begin{tabularx}{\linewidth}{>{\bfseries}p{0.12\linewidth}X>{\raggedright\arraybackslash}p{0.18\linewidth}}
\toprule
\textbf{ID} & \textbf{Current requirement and reason} & \textbf{Status} \\
\midrule
REQ-RS-01 & Complete a 2 m remotely commanded path with the operator outside the room; demonstrates remote entry. & Partly supported: $>2$ m user-reported with rover PoC; measured record and operator position open. \\
REQ-RS-02 & Display live phone video with latency at most 1 s; supports obstacle response. & Partly supported: laptop/desktop-to-phone stream reported; Pi/rover path and latency open. \\
REQ-RS-03 & Map at least three tagged temperature, humidity, and relative-gas records; demonstrates spatial variation. & Partly supported: synthetic schema/map passed; physical readings and rover map open. \\
REQ-RS-04 & Flag a potential anomaly at 125\% or more of baseline; this is a relative indicator, not a safety limit. & Partly supported: synthetic 124/125 boundary passed; sensor and integrated alarm open. \\
REQ-RS-05 & Map walls within $\pm$10 cm in at least four of five runs; preserves the proposal accuracy target. & Partly supported: ToF live map and synthetic scoring; eight-point tape comparison across five runs open. \\
REQ-RS-06 & Deliver the completed summary to the watch within 5 s; provides portable results. & Partly supported: phone photo transfer in SquareLine3 reported; rover summary and timing open. \\
REQ-RS-07 & Keep total hardware cost at or below CAD \$700; stakeholder constraint N7. & Open: CAD \$675 estimate is provisional; final cost proof open. \\
REQ-RS-08 & Build firmware from a clean copy using one documented command; supports reproducibility. & Partly supported: PoC builds passed separately; full integrated clean-copy build open. \\
REQ-RS-09 & Use the exact Waveshare watch to display the map, summary, and anomaly indications. & Partly supported: Wi-Fi/UI and phone-uploaded image display reported; final survey content open. \\
REQ-RS-10 & Include a custom PCB in the final system; breadboard/modules alone are insufficient for the final build. & Open: binding plan; PCB not yet designed or built. \\
\bottomrule
\end{tabularx}

\newpage
\section*{5. Risk Register}

Priority estimates use L-09's 1--4 probability ($p$) and consequence ($c$) indices; $P=p\times c$. These are provisional judgments, not observed failure frequencies. R1--R3 continue the proposal's IDs; R4--R5 arose during controller and watch PoCs.

\begin{tabularx}{\linewidth}{>{\bfseries}p{0.07\linewidth}Xcccc}
\toprule
ID & \textbf{Risk} & $p$ & $c$ & $P$ & \textbf{Rank} \\
\midrule
R1 & Wall error from slip/range noise & 3 & 4 & 12 & 1= \\
R3 & Video or watch delivery delayed by wireless traffic & 3 & 4 & 12 & 1= \\
R2 & Relative-gas baseline drift or slow response & 3 & 3 & 9 & 3 \\
R4 & Controller load, timing, or power exceeds budget & 2 & 3 & 6 & 4 \\
R5 & Final watch UI exceeds display/app limits & 2 & 2 & 4 & 5 \\
\bottomrule
\end{tabularx}

\begin{center}
\textbf{Risk matrix: $P=p\times c$ (IDs show estimated positions)}\\[-2pt]
\begin{tabular}{c|cccc}
 & $c=1$ & $c=2$ & $c=3$ & $c=4$ \\
\hline
$p=4$ & 4 & 8 & 12 & 16 \\
$p=3$ & 3 & 6 & 9 (R2) & 12 (R1, R3) \\
$p=2$ & 2 & 4 (R5) & 6 (R4) & 8 \\
$p=1$ & 1 & 2 & 3 & 4
\end{tabular}
\end{center}

\textbf{R1, tied highest.} \textbf{Statement:} slip/range noise may exceed $\pm$10 cm. \textbf{Assumption:} odometry and ranging are repeatable. \textbf{Consequence:} REQ-RS-05 fails and mapped locations mislead. \textbf{Probability:} $p=3$, judged likely because no physical accuracy run exists. \textbf{Evidence:} planned: eight tape-measured wall points in each of five runs. \textbf{Mitigation:} calibrate geometry, slow the rover, combine encoder/range data, or revise the mapping method.

\textbf{R3, tied highest.} \textbf{Statement:} concurrent traffic may delay phone video or watch summary. \textbf{Assumption:} both links meet the 1 s and 5 s limits together. \textbf{Consequence:} REQ-RS-02/06 fail. \textbf{Probability:} $p=3$, judged likely while only separate PoCs exist. \textbf{Evidence:} planned: timed concurrent doorway video and summary trials. \textbf{Mitigation:} reduce data rate, buffer/retry, or test another link.

R2 needs safe baseline, humidity, dwell-time, and repeatability tests; change sampling or sensor if unstable. R4 needs load/power measurements; divide control if needed. R5 needs the final survey UI build and flash test; V2 used 93\% of its selected app partition, while V3 used another partition.

\section*{6. Gate 1 Readiness}

The evidence supports a \textbf{preliminary proceed-with-conditions} position. Physical V1/V2 watch results (L-06), user-reported photo display in SquareLine3 (L-10), phone drive with a ToF-based live map, and the separate laptop-to-phone video record show useful subsystems. Ten synthetic tests support the REQ-RS-03--05 software logic. None passes the full physical acceptance conditions: operator position and calibrated path, Pi/rover camera latency, real tagged environmental readings, five tape-measured mapping runs, final watch map/summary/anomaly, and 5 s delivery remain open. All \texttt{src/} code is PoC. The final custom PCB and full integrated clean-copy build also remain open.

Before the next gate, document the physical acceptance trials, test the selected BME680/BME688 baseline, measure the power budget, and build the integrated system from a clean copy. Revisit the L-09 concept scores and risk estimates after those tests. These gaps prevent a claim of complete feasibility today.

\begin{thebibliography}{9}
\small
\bibitem{waveshare}
Waveshare, ``ESP32-S3-Touch-AMOLED-2.06 Product Page and Wiki,'' accessed October 23, 2026.

\bibitem{espressif}
Espressif Systems, ``ESP32-S3 Series Datasheet'' and ``RF Coexistence Guide,'' accessed October 23, 2026.

\bibitem{bosch}
Bosch Sensortec, ``BME688 Gas Sensor Product Information,'' accessed October 23, 2026.
\end{thebibliography}

\end{document}
